\chapter{MicroTE: Fine Grained Traffic Engineering for Data Centers}
\paperinfo{T. Benson, A. Anand, A. Akella (U. Wisconsin--Madison), M. Zhang (Microsoft Research)}{ACM CoNEXT 2011}

\begin{keybox}[The paper in three sentences]
Data center traffic looks random over minutes, but over \key{1--2 seconds} a good part of it is \key{predictable}: a rack pair
that sent about 50\,MB in the last second will send about 50\,MB in the next one. MicroTE uses a central OpenFlow controller to
find these predictable rack pairs every second and \key{route them carefully first} (minimising the most-loaded link). The
remaining unpredictable traffic is spread with \key{weighted ECMP} over the capacity that is left. When everything is
predictable MicroTE acts like optimal routing; when nothing is, it falls back to ECMP.
\end{keybox}

% ------------------------------------------------------------------
\section{The problem: how should traffic be spread inside a data center?}

A data center network has many paths between any two racks. In a two-tier tree with 2 core switches, a packet from rack A to
rack B can go up through core 1 or core 2; in a fat-tree there are many more choices. \key{Traffic engineering (TE)} is the job
of deciding which traffic goes on which path so that no link is overloaded.

The metric the paper uses throughout is the \key{Maximum Link Utilisation (MLU)}: the load on the busiest link divided by its
capacity. If MLU $> 1$, that link is receiving more than it can carry, so packets are queued and then dropped. A good TE scheme
keeps MLU low every second.

\paperfig[0.5]{mt_fig1}{MLU each second over a 2-hour trace for Spanning Tree, ECMP and optimal routing on a canonical tree.
Anything to the right of 1.0 means loss. ECMP is much better than a single tree, but still 15--20\% behind optimal.\label{fig:mtmlu}}

\subsection{What data centers did at the time, and why each falls short}

The authors replayed real traces from two data centers (CLD: a 1500-server cloud cluster running MapReduce jobs; UNV: a university
data center running 3-tier web services) through a simulator and compared each scheme against an \key{optimal} router. The optimal
router is an oracle: it knows the exact traffic of the next second and solves a linear program to minimise the MLU.

\begin{center}
\small
\begin{tabularx}{\linewidth}{@{}lXX@{}}
\toprule
Scheme & How it routes & Why it is not good enough\\
\midrule
\key{Spanning Tree} & All traffic follows one loop-free tree. & Most links are never used, so the tree links overload.\\
\key{ECMP} & Each flow is hashed onto one of the equal-cost paths, at random. & Balances the \emph{number} of flows, not their \emph{bytes}. Two
bursty rack pairs can be hashed onto the same link and burst together. It never looks at load, so it never fixes this.\\
\key{Fat-tree flow scheduling} & Every few seconds, each lower switch looks at its own uplinks and moves a few big flows from a busy port to a quiet one. &
The decision is \emph{local}. Two switches can each move a flow onto the same core switch, which then sends both down the same link. 23\% from optimal.\\
\key{VL2} & Valiant load balancing, implemented with ECMP per flow. & Inherits ECMP's blindness to load; up to 20\% from optimal.\\
\key{Hedera} & ECMP for small flows; a central scheduler places flows larger than 100\,MB. & In real traces, the traffic that collides is mostly in flows
\emph{under} 100\,MB, so Hedera behaves like ECMP.\\
\bottomrule
\end{tabularx}
\end{center}

\textbf{A small example of why ECMP fails.} Rack pairs (A$\to$B) and (C$\to$D) each send only 200\,Mb/s \emph{on average}, but in
bursts of 800\,Mb/s. ECMP hashes both onto the path through core 1. When both burst in the same second, that 1\,Gb/s link needs
1.6\,Gb/s, so MLU $= 1.6$ and packets are lost, even though the path through core 2 is idle. ECMP will never notice, because it
never measures load and never moves a flow once it has hashed it.

\begin{keybox}[Three reasons existing TE fails]
\begin{enumerate}
  \item \textbf{No multipath} (Spanning Tree): capacity is wasted.
  \item \textbf{Not load-aware} (ECMP, VL2): flows are spread by count, not by how many bytes they actually send right now.
  \item \textbf{No global view} (fat-tree's local heuristic): locally sensible moves collide somewhere else.
\end{enumerate}
\end{keybox}

\subsection{Why not use ISP-style traffic engineering?}
ISPs do TE with a central planner that predicts the traffic matrix and computes routes, but they work on time scales of
\emph{hours} and assume traffic is smooth. Data center traffic is very bursty: earlier measurements showed 80\% of flows last less
than 10\,s, about 100 new flows arrive every millisecond, and the traffic matrix is \emph{not} predictable over 100\,s or more. A
route computed for the average of the last hour is useless for the next second. So the question is: is there \emph{any} time scale at
which data center traffic is predictable?

% ------------------------------------------------------------------
\section{The key observation: short-term predictability}

The authors look at the traffic matrix between \key{ToR (top-of-rack) switch pairs}, one entry per (source rack, destination rack),
measured every second. They call an entry \key{predictable} if it changes by less than 20\% from one second to the next.

\paperfig[0.85]{mt_fig2}{(a) For each 1\,s interval, what fraction of all traffic came from ToR pairs whose volume stayed within 20\%. (b) For
each ToR pair, how long (on average) it stays stable.\label{fig:mtpred}}

What they found (Figure~\ref{fig:mtpred}):
\begin{itemize}
  \item \textbf{A large part of the traffic is predictable for the next second.} In CLD, in 80\% of the 1\,s intervals, more than
        27\% of all bytes came from ToR pairs that stayed stable. In UNV, in over 65\% of intervals, \emph{all} the traffic was stable.
  \item \textbf{It stays predictable for about 1.5--2.5\,s.} In CLD, 60\% of ToR pairs stay stable for 1.6--2.5\,s on average; in UNV,
        35\% of pairs stay stable for 1.5--5\,s.
\end{itemize}

\textbf{Why traffic behaves like this (the applications explain it).}
\begin{itemize}
  \item \textbf{CLD (MapReduce).} The master sends small task commands to workers (little traffic), then workers send their results
        back (lots of traffic). The scheduler tries to put workers in the same rack as the master, but big jobs spill over to other racks.
        While those remote workers stream results back, the rack-to-rack volume is steady for a couple of seconds: that is the
        predictable part. Control messages and asynchronous block replication are the unpredictable part.
  \item \textbf{UNV (3-tier web).} Each rack holds one tier (all web front-ends in one rack, all business-logic servers in another). The
        front-end only talks to the business tier, which only talks to the database tier, so the same rack pairs keep talking. But the
        messages are small and depend on user requests, so stability is short.
\end{itemize}

\begin{keybox}[Design requirements derived from the study]
\begin{enumerate}
  \item \textbf{Use multiple paths.}
  \item \textbf{Coordinate routing with a global view} of traffic, not local decisions.
  \item \textbf{Exploit short-term predictability and adapt within about a second}; and when traffic is \emph{not} predictable, do
        no worse than ECMP (otherwise no operator will adopt it).
\end{enumerate}
\end{keybox}

% ------------------------------------------------------------------
\section{MicroTE's idea in one rule}

\begin{keybox}[The one rule]
Every second: \textbf{route the predictable ToR pairs optimally first}, then \textbf{spread the unpredictable traffic with weighted ECMP
over whatever capacity is left}.
\end{keybox}

Why this split makes sense: for a predictable pair we know roughly how many bytes it will send in the next second, so we can place it
precisely, like a puzzle piece. For an unpredictable pair we do not know its next-second volume, so careful placement would be based on a
wrong number. The best we can do is spread it out, but we can at least avoid the links that the predictable traffic has already filled.

\textbf{It degrades gracefully}:
\begin{center}
\begin{tabular}{@{}lll@{}}
\toprule
Fraction of traffic predictable & What MicroTE does & Behaves like\\
\midrule
almost all & routes nearly everything with full knowledge & optimal routing\\
some & routes the known part carefully, spreads the rest around it & between ECMP and optimal\\
none & nothing to route carefully; weights are equal & plain ECMP\\
\bottomrule
\end{tabular}
\end{center}

% ------------------------------------------------------------------
\section{Architecture: three components}

\paperfig[0.55]{mt_fig3}{MicroTE architecture. Servers (monitoring component) send flow statistics to the network controller; the routing
component turns the global view into routes; the controller installs them in the switches.\label{fig:mtarch}}

\begin{enumerate}
  \item \key{Monitoring component}: a kernel module on every server. It measures how much the server sends to whom, and one server per rack
        summarises this for the whole rack.
  \item \key{Network controller}: an OpenFlow (NOX) controller. It collects the rack summaries into one global traffic matrix, learns the
        topology, and installs flow entries into the switches.
  \item \key{Routing component}: takes the global view from the controller and computes the routes, using an LP or a fast bin-packing heuristic.
\end{enumerate}

\subsection{Monitoring: why measure at the servers, not the switches?}
Two options exist. (1) The controller polls every switch's flow counters via OpenFlow. (2) The servers measure their own traffic and report it.
MicroTE chooses the servers, for three reasons:
\begin{itemize}
  \item \textbf{Proactive, triggered updates.} A server can report the moment its traffic changes. OpenFlow switches at the time could only be
        \emph{polled}, so the controller would have to ask every switch every second.
  \item \textbf{Less control traffic.} Polling every switch every second creates a lot of messages in the network.
  \item \textbf{Moves the work off the switches.} Counting and summarising flows uses the server's CPU (under 4\%) instead of the switch's weak
        control processor.
\end{itemize}

\subsection{What the monitoring component does, step by step}
\begin{enumerate}
  \item \textbf{Bootstrapping.} When a new server plugs into a ToR switch, the switch tells the controller (a normal OpenFlow event). The
        controller now knows which rack the server is in. The \emph{first} server to appear in a rack is made the \key{designated server}
        for that rack. The controller tells each new server where the controller is and who its rack's designated server is.
  \item \textbf{Measure.} Every server's kernel module records its server-to-server traffic every \key{0.1\,s}, and once per second sends these
        readings to its rack's designated server.
  \item \textbf{Summarise.} The designated server knows the server-to-rack mapping (sent by the controller). It adds up all server-to-server
        bytes into \key{ToR-to-ToR} volumes. This turns a huge matrix into a small one (75 racks instead of 1500 servers in CLD).
  \item \textbf{Decide predictability.} For each ToR pair, it keeps a running \emph{mean} of recent volumes. If the current volume is within
        \key{$\delta$ = 20\%} of the mean, the pair is marked \key{predictable}, and its mean is used as the forecast for the next second.
  \item \textbf{Report only changes.} Instead of resending the whole matrix every second, it sends an update only when the \emph{set} of
        predictable pairs changes, or when a predictable pair's volume moves by more than $\delta$.
\end{enumerate}

\textbf{Example of the predictability test} ($\delta = 20\%$), for one ToR pair over a few seconds:
\begin{center}
\begin{tabular}{@{}ccccl@{}}
\toprule
Second & Volume (MB) & Running mean & Within 20\%? & Status\\
\midrule
1 & 50 & 50 & -- & first sample\\
2 & 54 & 52 & $|54-52|/52 = 4\%$ \checkmark & predictable, forecast 52\,MB\\
3 & 48 & 50.7 & 5\% \checkmark & predictable\\
4 & 120 & 68 & 76\% $\times$ & unpredictable $\to$ goes to weighted ECMP\\
\bottomrule
\end{tabular}
\end{center}
Second 4 also changes the set of predictable pairs, so this triggers an update to the controller.

\subsection{Network controller}
\begin{itemize}
  \item \textbf{Global traffic matrix}: built by joining the rack summaries.
  \item \textbf{Topology}: learned from OpenFlow's link/switch up-and-down messages.
  \item \textbf{When to recompute}: on any update (new set of predictable pairs, or a predictable pair changing by more than $\delta$), but
        \key{at most once per second}. This prevents a burst of updates from causing a burst of recomputations.
  \item \textbf{Install routes}: sends the routing component's paths to the switches as OpenFlow flow entries.
\end{itemize}

% ------------------------------------------------------------------
\section{The routing component}

\paperfig[0.48]{mt_fig4}{What the routing component does with each new global view.\label{fig:mtflow}}

Following Figure~\ref{fig:mtflow}, for each new global view:
\begin{enumerate}
  \item Separate the ToR pairs into predictable and unpredictable.
  \item \textbf{Route the predictable pairs} (with the LP or the bin-packing heuristic below), using their forecast volumes.
  \item \textbf{Weighted ECMP for the rest.} For each unpredictable pair, its flows are spread over the equal-cost paths, with each path's
        weight proportional to the capacity \emph{left over} after the predictable traffic is placed.
  \item Save the view in the history (used for the running means).
  \item If the new routes differ significantly from the installed ones, return them for installation; otherwise do nothing (avoids
        needless churn in the switches).
\end{enumerate}

\subsection{Option 1: Linear program (exact)}
Write $T_{u,v}$ for the forecast volume of ToR pair $(u,v)$, $P_{u,v,k}$ for the links on its $k$-th equal-length path ($1\le k\le K$),
and $f_{u,v,k}$ for the fraction of the pair's traffic put on path $k$.
\[
\text{minimise } util \quad\text{subject to}\quad
\sum_{k=1}^{K} f_{u,v,k} = 1 \;\;\forall (u,v), \qquad
\sum_{(u,v,k):\, e \in P_{u,v,k}} T_{u,v}\, f_{u,v,k} \;\le\; util \cdot Cap_e \;\;\forall e.
\]
In words: all of each pair's traffic must be placed on its $K$ paths, and every link's load must be at most $util$ times its capacity.
Minimising $util$ minimises the MLU. The LP may split a pair over several paths.

\textbf{Problem}: it is slow as the network grows. With 100 ToRs and 50\% of pairs predictable it takes 0.6\,s; with 200 ToRs it takes 0.5\,s
even at 5\% predictable, and 28\,s at 25\%. MicroTE must finish within its 1\,s cycle, so the LP is only usable in small data centers.

\subsection{Option 2: Bin-packing heuristic (fast)}
\begin{enumerate}
  \item Sort predictable ToR pairs by forecast volume, \textbf{largest first}.
  \item For the next pair, look at its $K$ equal-hop paths. The \key{cost of a link} is $1/(\text{remaining capacity})$; a path's cost is the
        sum over its links. Pick the \textbf{cheapest} path.
  \item Put the \emph{whole} pair on that path and subtract its volume from the remaining capacity of each link on the path (which makes those
        links more expensive for later pairs).
  \item Repeat for the next pair.
\end{enumerate}
Big pairs go first because they are the hardest to fit; links that are already full become expensive, so later pairs avoid them.

\textbf{Worked example.} Two ToR pairs from rack A to rack B. There are two paths: via core 1 (links A--C1, C1--B) and via core 2
(links A--C2, C2--B), every link 10\,Gb/s. The predictable pairs (forecast for the next second) are P1 = 6\,Gb/s and P2 = 5\,Gb/s,
and there is also an unpredictable pair U of about 3\,Gb/s.
\begin{center}
\small
\begin{tabular}{@{}llll@{}}
\toprule
Step & Path costs ($\sum 1/\text{remaining}$) & Choice & Remaining capacity after\\
\midrule
P1 (6, largest) & core1: $1/10+1/10 = 0.2$; core2: $0.2$ & core 1 (tie) & core1 links: 4; core2 links: 10\\
P2 (5) & core1: $1/4+1/4 = 0.5$; core2: $0.2$ & core 2 & core1: 4; core2: 5\\
U (unpredictable) & weights $\propto$ leftover: core1 4, core2 5 & split U's flows 4:5 & core1 $\approx$ 1.3 used; core2 $\approx$ 1.7 used\\
\bottomrule
\end{tabular}
\end{center}
Final link loads: core 1 path about 7.3\,Gb/s, core 2 path about 6.7\,Gb/s, so MLU $\approx 0.73$.

\textbf{Same input, plain ECMP}: each pair is hashed to a path at random. With probability $1/2$ P1 and P2 land on the same path:
$6+5 = 11$\,Gb/s on a 10\,Gb/s link, so MLU $\ge 1.1$ and there is loss, while the other path carries only U.

\textbf{Running time}: one shortest-path search per pair, $O(E + N\log N)$, so $O(PN\log N + PE + P\log P)$ for $P$ predictable pairs
(the last term is the sort). For 200 ToRs and 50\% predictable pairs it takes under 0.25\,s. The price is that it never splits a pair over
paths, so it is a bit less optimal than the LP.

\textbf{Which to use}: LP for small data centers (about 50 ToRs, or 100 ToRs with few predictable pairs); bin-packing for medium and large ones.

% ------------------------------------------------------------------
\section{Implementation and overheads}
\begin{itemize}
  \item \textbf{Code}: the controller and routing component are C++ modules in NOX; the monitor is a Linux kernel module in C. The switches only need
        OpenFlow firmware, and the servers need the kernel patch. No new hardware or topology.
  \item \textbf{Demand reports}: 8 bytes per communicating ToR pair (2\,B source, 2\,B destination, 4\,B volume). Because the traffic matrix is sparse,
        this is about 4\,MB/s even for 10K servers, under 0.5\% of a server's uplink.
  \item \textbf{Route installation}: 100\,B per OpenFlow entry, about 50\,MB in total for 10K servers (under 5\% of a link). Installing routes took
        under 40\,ms with 10 flows per server and under 150\,ms with 80 flows per server, which is fast enough for a 1\,s cycle.
\end{itemize}

% ------------------------------------------------------------------
\section{Results in brief}
\begin{itemize}
  \item At ToR-pair granularity MicroTE is always better than ECMP, but 1--15\% from optimal. The reason: predictability really comes from
        individual applications and servers, and grouping them into racks blurs it.
  \item At \key{server-to-server} granularity, MicroTE gets within \key{1--5\%} of optimal and cuts lost bytes by 65\% compared with ECMP
        under heavy load. The cost is more control traffic and computation.
  \item When predictability is low, MicroTE behaves like ECMP; when it is high (and traffic is heavy), ECMP's MLU goes past 1 while MicroTE stays
        close to optimal.
\end{itemize}

% ------------------------------------------------------------------
\section{Summary comparison}
\begin{center}
\small
\begin{tabularx}{\linewidth}{@{}lXXXX@{}}
\toprule
 & ECMP / VL2 & Fat-tree scheduler & Hedera & \textbf{MicroTE}\\
\midrule
Multipath & yes & yes & yes & yes\\
Load-aware & no & yes, local & only flows $>$100\,MB & \key{yes, per ToR pair}\\
Global view & no & no & yes for big flows & \key{yes}\\
Time scale & per flow, fixed & every few s & every few s & \key{every 1\,s}\\
When traffic is unpredictable & -- & -- & -- & falls back to ECMP\\
\bottomrule
\end{tabularx}
\end{center}

\begin{keybox}[Exam one-liner]
MicroTE is OpenFlow-based traffic engineering for data centers. It relies on the fact that a good part of rack-to-rack traffic is predictable
for about 1--2\,s. Servers measure their traffic every 0.1\,s; a designated server per rack summarises it into ToR-to-ToR volumes and marks a pair
predictable if it is within 20\% of its running mean. Every second the controller routes the predictable pairs first (LP for small DCs, greedy
bin-packing with link cost $1/\text{remaining capacity}$ for large ones), then spreads the rest with ECMP weighted by leftover capacity. Result:
close to optimal when traffic is predictable, ECMP when it is not.
\end{keybox}

\section{Revision checklist}
\begin{itemize}
  \item MLU = busiest link's load / capacity; $>1$ means loss. Existing schemes are 15--23\% worse than optimal.
  \item Three failures: no multipath (spanning tree), not load-aware (ECMP, VL2), no global view (fat-tree local moves); Hedera's 100\,MB threshold misses the colliding flows.
  \item ISP TE works over hours; DC traffic is unpredictable over 100\,s or more, but predictable over 1--2\,s (MapReduce result transfers, tiered web apps).
  \item Rule: predictable pairs routed optimally first, unpredictable via weighted ECMP on leftover capacity.
  \item Components: server monitors (0.1\,s readings, designated server per rack, ToR aggregation, 20\% test, send only changes), NOX controller (global TM, topology, recompute at most once per second), routing component.
  \item Why measure at servers: triggered updates, less polling traffic, off-loads switches.
  \item LP minimises $util$; exact but slow. Bin-packing: largest first, link cost $1/\text{remaining}$, cheapest of $K$ paths, no splitting; fast.
  \item Server-to-server granularity brings MicroTE within 1--5\% of optimal.
\end{itemize}
